\documentclass{article}
\usepackage[utf8]{inputenc}
\usepackage[T1]{fontenc}
\usepackage{geometry}
\usepackage{graphicx}
\usepackage{float}
\usepackage{hyperref}
\usepackage[english,indonesian]{babel}

\title{Optimasi Pemrosesan Template LaTeX Menggunakan Pendekatan Berbasis AI}
\author{Academic Support Agent}
\date{\today}

\begin{document}
\maketitle

\begin{abstract}

\end{abstract}

\chapter{PENDAHULUAN}


\chapter{TINJAUAN PUSTAKA}


% Original content (commented out):
% \documentclass[12pt]{article}
\linespread{1.2}
\usepackage[utf8]{inputenc}
\usepackage[english, swedish]{babel}
\usepackage{graphicx}
\usepackage{biblatex}
\addbibresource{bibliography.bib}
\usepackage[a4paper,pdftex,bottom=20mm, width=160mm]{geometry} % A4paper margins

\usepackage{fancyhdr}

\pagestyle{fancy}
\fancyhf{} % Rensa sidhuvud och sidfot
\fancyhead[L]{Realgymnasiet}
\fancyhead[R]{Titel på gymnasiearbete}
\fancyfoot[C]{\thepage} % Sidnumrering i sidfoten
\renewcommand{\headrulewidth}{0pt} % Ta bort linjen från sidhuvudet

\title{
\includegraphics[scale=0.5]{realgymnasiet_svart_logo.png} \\
\vspace{2.0cm} \textbf{Titel på rapporten} \\
 \endgraf\rule{\textwidth}{.2pt}

   }
   
   
\author{Ditt namn}
\date{\today}

\begin{document}

\maketitle
\thispagestyle{empty}

\newpage
\thispagestyle{empty}

\begin{abstract}
\noindent
Ett abstract är en kort sammanfattning av hela ditt arbete. Det skriver du allra sist när du är klar med resten. Det bör inte vara mer än 100-150 ord. Se till att alla delar av arbetet är med. Ett exempel hittar du nedan: 

\begin{quote}

This study explored the pattern of video game usage and video game addiction among male college students and examined how video game addiction was related to expectations of college engagement, college grade point average (GPA), and on-campus drug and alcohol violations. Participants were 477 male, first year students at a liberal arts college. In the week before the start of classes, participants were given two surveys: one of expected college engagement, and the second of video game usage, including a measure of video game addiction. Results suggested that video game addiction is (a) negatively correlated with expected college engagement, (b) negatively correlated with college GPA, even when controlling for high school GPA, and (c) negatively correlated with drug and alcohol violations that occurred during the first year in college. Results are discussed in terms of implications for male students’ engagement and success in college, and in terms of the construct validity of video game addiction.


\end{quote}

\noindent
\textbf{Keywords}: Här skriver du 3-5 ord som sammanfattar din uppsats. Oftast är det substantiv som beskriver innehållet, men även metod kan vara ett keyword. \\


\noindent
Exempel nedan: 
Keyword: Elbil, LIDAR, RADAR, Intervjustudie, Trafiksäkerhet

\end{abstract}

\newpage
\thispagestyle{empty}

\tableofcontents
\newpage
\section{Inledning}

I det här avsnittet ska du skriva en intresseväckande inledning som får läsaren att vilja fortsätta läsa din uppsats. Ett vanligt grepp är att ha en populärvetenskaplig ansats, det vill säga att inledning är lite mindre formell än resten av uppsatsen. Det är dock fortfarande viktigt att källhänvisa och skriva korrekt. 
Det är inte ovanligt att skriva om inledningen ett antal gånger. Börjar du ditt uppsatsarbete med att skriva inledning får du förmodligen skriva om den flera gånger. Det kan därför vara smart att avvakta lite tills du vet mer vilken riktning ditt arbete tar innan du skriver inledningen.


\subsection{Syfte}

Det här är en underrubrik till inledningen. Här ska du förklara varför du skriver din uppsats. Vad fyller den för (forsknings)syfte? Ditt syfte är aldrig ”att få betyg” eller ”klara kursen” utan det måste vara något mer genuint. Syftesförklaringen ska vara kort och formulerad som en eller flera ”att-satser”. Se exempel: 

Syftet med den här uppsatsen är

\begin{itemize}
    \item Att undersöka skillnaden i läshastighet mellan tryckta och digitala läromedel
    \item Att undersöka skillnader i kostnader och logistik mellan tryckta och digitala läromedel
\end{itemize}

\subsection{Mål}

Under den här rubriken skriver du målet eller målen med ditt arbete. Ska du bygga ett spel i Unity? Ska du kunna döda en boss? Hoppa över ett hål? Beskriv dina mål här: 

\begin{itemize}
    \item I spelet ska man kunna tjäna XP och gå upp i level
    \item I spelet ska man kunna samla resurser och bygga egna verktyg

Tänk på att dina mål och dina frågeställningar ska hänga ihop. De kan ofta vara väldigt lika, det är inte ett problem. 

\end{itemize}


\subsection{Disposition}

Det här är en rubrik som ibland inte behövs, men det kan vara bra att förtydliga vissa saker kring uppsatsens disposition, alltså upplägg. 

Under den här rubriken kan du också förtydliga om det är något i din rapport som bryter mot normerna. Ska du till exempel använda kapitäler för att markera ut något eller har du döpt om ditt analyskapitel? Berätta det här då.

\newpage
\section{Bakgrund / Teori}

I det här kapitlet ska du ta upp sådant som sagts i ämnet före dig. Vad grundar du din kunskap på? Observera att det ju inte går att ta upp allt som sagts i ämnet så det gäller att välja vad som ska få plats. Ett bra sätt att tänka kan vara att ha en historisk ansats, det vill säga att först ta upp gamla källor, sedan lite nyare och avsluta med det allra nyaste. På så vis får läsaren en guidning inte bara genom den teoretiska bakgrunden utan också hur det har utvecklats. 

Om du gör ett praktiskt arbete så kan du under den här rubriken presentera lite andra projekt med liknande tillvägagångssätt som ditt. Om du jobbar i Unity kanske du kan presentera andra spel som är gjorda i det programmet. 


\newpage
\section{Metod och material}

Under den här övergripande rubriken ska du göra tre konkreta och olika saker. De förklaras i tre underrubriker. 


\subsection{Metod}

Under den här rubriken ska du berätta om de verktyg du använder för att uppfylla ditt praktiska arbete. Vilka programmeringsspråk ska du använda? Beskriv de kortfattat.


\subsection{Material}

Under den här rubriken beskriver du information du samlat in för att uppfylla ditt mål. Har du tittat på guider och genomgångar? Har du läst en massa text om ditt programmeringsspråk? Har du gjort användarundersökningar?


\subsection{Metoddiskussion}

Under den här rubriken finns det plats för en kort diskussion om din valda metod. Vad hade du kunnat göra istället? Observera att den här rubriken kan smälta samman med den övergripande Diskussionen. 



\newpage
\section{Genomförande}

Ibland heter den här rubriken ”Analys” eller ”Analys \& resultat”. Oavsett namn på rubriken ska du här beskriva hur du uppnått ditt eller dina mål. Det kan därför vara bra att ha en underrubrik för varje mål. Tänk på att din text här måste målen. Vill du ta upp något utöver det får du helt enkelt skriva till nya mål. 

Här passar det också med bilder från eller på din produkt. Beskriv dina problem och hur du löste dem. Glöm inte att förklara dina bilder och figurer. 



\newpage
\section{Diskussion}

Under den här rubriken ska du diskutera resultaten från föregående rubrik? Var de uppseendeväckande? Hade du kunnat göra på något annat sätt för att få ett annat resultat? Hur kan resultatet kopplas till ditt teorikapitel? 

Avsluta gärna med att ge förslag på hur du eller någon annan skulle kunna gå vidare med ditt arbete. Vad kan man göra mer?



\newpage
\section{Förklaring referenslista}

Vi använder oss av Harvard referenslista. Bra exempel hittar du på Umeås universitet, Harvard - skriva referenslista. (2019) https://www.umu.se/bibliotek/soka-skriva-studera/skriva-referenser/harvard-skriva-referenslista/) (Hämtad 2024-02-09). \cite{Umeå}
\\
\\
Bryman (2008) \cite{Social} betonar vikten av att använda en mångfald av forskningsmetoder för att uppnå noggranna och tillförlitliga resultat.

\printbibliography


\end{document}


\end{document}